\section{Cài đặt Cụm Hadoop Fully-Distributed}

% ------------------------------------------------------------
\subsection{Kiến trúc cụm vật lý}

Cụm Hadoop fully-distributed của nhóm gồm một node Master và hai node Worker, tất
cả được kết nối qua mạng LAN nội bộ. Tôi (MSSV 23120099, Trưởng nhóm Leader A) phụ
trách node Master chạy Arch Linux. Các thành viên còn lại của nhóm phụ trách các
node Worker. Trong chế độ fully-distributed, mỗi node đảm nhận vai trò chuyên biệt:
Master chạy các tiến trình quản lý (NameNode, ResourceManager), còn Worker chạy các
tiến trình lưu trữ và tính toán (DataNode, NodeManager).

\begin{center}
\begin{tabular}{lll}
\toprule
\textbf{Node} & \textbf{Vai trò} & \textbf{Địa chỉ IP} \\
\midrule
master (23120099) & NameNode / ResourceManager & 100.95.242.89 \\
tuanho            & DataNode / NodeManager     & 100.81.238.8 \\
msi               & DataNode / NodeManager     & 100.101.144.37 \\
\bottomrule
\end{tabular}
\end{center}

% ------------------------------------------------------------
\subsection{Cấu hình mạng và SSH}

Để các node trong cụm có thể giao tiếp với nhau bằng tên thay vì địa chỉ IP, tôi
chỉnh sửa file \texttt{/etc/hosts} trên tất cả các node, thêm các dòng ánh xạ tên
sang IP tương ứng. Ví dụ nội dung file \texttt{/etc/hosts} trên node Master:

\begin{minted}{bash}
# Cấu hình /etc/hosts cho cụm Hadoop fully-distributed
100.95.242.89    master
100.81.238.8     tuanho
100.101.144.37   msi
\end{minted}

Tiếp theo, tôi cấu hình SSH không mật khẩu (passwordless SSH) từ Master đến tất cả
Worker để các script \texttt{start-dfs.sh} và \texttt{start-yarn.sh} có thể tự động
SSH vào từng Worker để khởi động daemon mà không cần nhập mật khẩu:

\begin{minted}{bash}
# Tạo cặp khóa SSH trên node Master
ssh-keygen -t rsa -b 4096 -f ~/.ssh/id_rsa -N ""

# Sao chép public key sang từng Worker (thay username và IP thực tế)
ssh-copy-id tuanho@tuanho
ssh-copy-id truc2005@msi

# Kiểm tra kết nối SSH không mật khẩu
ssh tuanho@tuanho "hostname"
ssh truc2005@msi "hostname"
\end{minted}

Nội dung file \texttt{\$HADOOP\_HOME/etc/hadoop/workers} cần được cập nhật để liệt
kê tên hoặc IP của các Worker node:

\begin{minted}{bash}
tuanho
msi
\end{minted}

Đối với \texttt{core-site.xml}, giá trị của \texttt{fs.defaultFS} cần được đổi từ
\texttt{hdfs://localhost:9000} sang \texttt{hdfs://master:9000} để các Worker node
có thể kết nối đến NameNode trên Master thông qua tên hostname:

\begin{minted}{xml}
<configuration>
  <property>
    <name>fs.defaultFS</name>
    <value>hdfs://master:9000</value>
  </property>
</configuration>
\end{minted}

% TODO: Thêm ảnh chụp màn hình kiểm tra SSH
\begin{figure}[H]
  \centering
  \includegraphics[width=0.9\textwidth]{2-node.png}
  \caption{Kiểm tra kết nối SSH không mật khẩu từ Master đến Worker}
  \label{fig:ssh-test}
\end{figure}

% ------------------------------------------------------------
\subsection{Khởi động và kiểm tra YARN}

Sau khi cấu hình xong, tôi khởi động toàn bộ cụm từ node Master bằng hai script
quen thuộc. Khác với chế độ pseudo-distributed, lần này các daemon sẽ được khởi
động trên nhiều máy vật lý thực sự:

\begin{minted}{bash}
# Khởi động từ node Master
start-dfs.sh    # Tự động SSH vào Worker để khởi động DataNode
start-yarn.sh   # Tự động SSH vào Worker để khởi động NodeManager
\end{minted}

Sau khi khởi động, tôi kiểm tra trạng thái trên từng node bằng \texttt{jps}:

\begin{minted}{bash}
# Trên Master — kỳ vọng có: NameNode, ResourceManager, SecondaryNameNode
jps

# Trên mỗi Worker — kỳ vọng có: DataNode, NodeManager
ssh tuanho@tuanho "jps"
ssh truc2005@msi "jps"
\end{minted}

Giao diện WebUI của NameNode có thể truy cập tại
\texttt{http://<master-ip>:9870} để xem trạng thái HDFS và danh sách DataNode. Giao
diện YARN ResourceManager có thể truy cập tại
\texttt{http://<master-ip>:8088} để theo dõi các job đang chạy và lịch sử thực thi.

% TODO: Thêm ảnh chụp màn hình WebUI NameNode
\begin{figure}[H]
  \centering
  \includegraphics[width=0.9\textwidth]{webui-namenode.png}
  \caption{WebUI NameNode hiển thị các DataNode trong cụm fully-distributed}
  \label{fig:webui-cluster}
\end{figure}

% TODO: Thêm ảnh chụp màn hình WebUI YARN
\begin{figure}[H]
  \centering
  \includegraphics[width=0.9\textwidth]{webui-yarn.png}
  \caption{WebUI YARN ResourceManager tại cổng 8088}
  \label{fig:webui-yarn}
\end{figure}
